\begin{center}
    {\large \textbf{Giorno 13} \textbar\ In viaggio con Ketut \textbar\ 21 Agosto 2024 \textbar\ \textbf{Ubud}, Bali}
\end{center}

\noindent\rule{\textwidth}{0.4pt}
\begin{figure}[h!] % Portrait images below
    \centering
    % First portrait image (on the left)
    \begin{minipage}[h]{0.48\textwidth} % Set width equal to half the page
        \centering
        \includegraphics[width=\textwidth]{images/flags/Screenshot templi.png}
    \end{minipage}
    \hfill
    % Text
    \begin{minipage}[h]{0.48\textwidth}
        \raggedright
        \textbf{Superficie}: 4.567 km² \\
        \vspace{0.3cm}
        \textbf{Popolazione}: 3,3 milioni \\
        \vspace{0.3cm}
        \textit{L'isola}
    \end{minipage}
\end{figure}

%%% TODO: CONTROLLARE FOTO DUPLICATE TRA INIZIO E FINE

\landscapeLargeMargin{images/day13/PXL_20240821_004821591.MP.jpg}

\landscapeLargeMargin{images/day13/PXL_20240821_005157529.MP.jpg}

\landscapeLargeMargin{images/day13/PXL_20240821_002908010.MP.jpg}

\landscapeLargeMargin{images/day13/PXL_20240821_004613667.jpg}

\landscapeLargeMargin{images/day13/PXL_20240821_033408219.MP.jpg}

\landscapeLargeMargin{images/day13/PXL_20240821_032929504.jpg}

\landscapeLargeMargin{images/day13/PXL_20240821_033139413.MP.jpg}

\landscapeLargeMargin{images/day13/PXL_20240821_054227343.jpg}

\landscapeLargeMargin{images/day13/PXL_20240821_052952611.MP.jpg}

\landscapeLargeMargin{images/day13/PXL_20240821_052612476.jpg}

\landscapeLargeMargin{images/day13/PXL_20240821_070418225.jpg}

\landscapeLargeMargin{images/day13/PXL_20240821_082006639.jpg}

\landscapeLargeMargin{images/day13/PXL_20240821_134507995.jpg}

\noindent 21 Agosto 2024 \newline

Che male la sveglia alle 5:30. Come nei migliori tempi del Messico, mi sono svegliato alle tre con la sensazione di aver dimenticato qualcosa. Mi è venuto il dubbio che la partenza non fosse alle sei bensì alle cinque, così durante la notte ho ricontrollato i messaggi con Ketut. Ovviamente era tutto nella mia immaginazione, ma fino alla sveglia sono rimasto in dormiveglia.

Quando l'allarme finalmente suona sono subito in piedi e in men che non si dica tiro giù dal letto anche Elisa. A cinque minuti dall'ora dell'appuntamento siamo fuori dalla stanza e alla reception sono già pronti con le breakfast box, in sacchetto brandizzato con il nome Gayatri scritto in oro. L'autista è puntualissimo, spacca il minuto, e quando saliamo in macchina ci salutiamo appena perché siamo mezzi addormentati — così non riesco a capire se sia Ketut o un suo collega. Si vede che siamo vicini alla fine del viaggio perché la mia voglia di socializzare è già diminuita. O forse è solo perché è ancora notte.

Usciti da Ubud mi sembra per un attimo di rivivere il Rally del Chiapas, con accelerate e frenate folli, salvo che il nostro autista è molto ligio e non supera i 50 km/h. Elisa crolla subito mentre io rimango un po' sveglio seguendo la strada su Maps e quando vedo che stiamo per raggiungere Goa Lawah chiedo all'autista se ci fermeremo subito lì o proseguiremo oltre. Mi conferma che quel tempio è ancora chiuso e che ci fermeremo al ritorno, quindi proseguiamo fino alla prima tappa reale: Tirta Gangga. Nell'ultimo tratto riesco a dormire un po' e quando arriviamo l'autista si vanta vittorioso di aver impiegato meno di due ore.

Come prima cosa decidiamo di fare colazione: ci sediamo per terra su un gradino, tra le bancarelle semi-aperte dei venditori disinteressati, con le nostre breakfast box. La colazione è di livello anche questa volta: ottimo pane dalla crosta croccante con cioccolata spalmabile e burro d'arachidi racchiusi in foglie di banana, accompagnato da frutta — banane, mandarini e salak. Divoriamo tutto con piacere e l'unica mancanza è una buona tazza di caffè, ma ci rifaremo presto.

Ci avviamo tra le bancarelle verso l'entrata del tempio. Qui vendono bottigliette di plastica piene di cibo per pesci, tutte da mezzo litro o da un litro e mezzo, con pallini colorati. Davanti all'entrata troviamo una bancarella dove è possibile fotografarsi tenendo in braccio animali selvatici: una grossa iguana, due giganteschi pipistrelli addomesticati appesi a testa in giù, un pitone albino e uno reale, e tre zibetti — i famosi mammiferi che mangiano ed espellono i chicchi di caffè. Eli scatta qualche foto mentre osserviamo una ragazza inglese andare fuori di testa per gli zibetti, che sembrano dei gattoni.

Facciamo i biglietti ed entriamo a Tirta Gangga — a mio modesto parere, il tempio più bello di Bali.

\begin{center}
    % Decorative line
    \noindent\rule{0.6\textwidth}{0.4pt}
\end{center}

\textit{Il nome Taman Tirta Gangga si traduce come "Acqua sacra del Gange". Il complesso fu costruito nel 1946 dall'ultimo re di Karangasem, I Gusti Bagus Jelantik. Tirta Gangga era destinato a essere un luogo di svago per il re e la sua famiglia. Fu distrutto quasi interamente dall'eruzione del vicino Monte Agung nel 1963 e successivamente ricostruito. Il complesso copre un ettaro e le sue acque vengono utilizzate per l'irrigazione, l'attività economica e la ricreazione.}

\begin{center}
    % Decorative line
    \noindent\rule{0.6\textwidth}{0.4pt}
\end{center}

L'interno è pieno di enormi piscine, attraversabili tramite passerelle costituite da colonnine di pietra a filo d'acqua, abitate da tantissime carpe koi. Capiamo subito il motivo del cibo per pesci in vendita all'ingresso: le carpe hanno dimensioni spropositate. L'afflusso quotidiano di turisti è altissimo e se ognuno compra una bottiglia di mangime da un litro e mezzo questi pesci rischiano di diventare obesi. Alcuni di loro hanno sviluppato muscoli della bocca innaturali, segno che passano l'intera giornata a mangiare. Mi fanno quasi pena — anche i pesci sono diventati un'attrazione per Instagram.

Siamo tra i primi a entrare e questo ci permette di scattare qualche foto senza altri turisti, in posa sulle colonnine a filo d'acqua, circondati dalle carpe. Il panorama circostante è da cartolina: trascorriamo un'ora abbondante a visitare le innumerevoli piscine, i gazebo adibiti alla preghiera ancora profumati d'incenso, a passeggiare accanto ai giardini in fiore e ad attraversare i ponti ornati di sculture di dragoni. Ci perdiamo in questo paradiso terrestre e con Elisa ci ritroviamo praticamente alla fine del percorso dopo esserci goduti in autonomia le bellezze del luogo.

Scopriamo che le antiche residenze reali del complesso sono ora state adibite a hotel, così che i turisti possano godere della vista delle piscine direttamente dalla loro stanza. A un prezzo ridicolo, vorrei aggiungere.

Ci vorremmo fermare ancora, ma ci aspetta una giornata piena, così torniamo verso la macchina dove ci attende l’autista. Scambiamo due parole e dalla sua difficoltà nel parlare inglese mi rendo conto che non si tratta di Ketut, che il giorno prima aveva dimostrato un'ottima conoscenza della lingua.

\begin{center}
    % Decorative line
    \noindent\rule{0.6\textwidth}{0.4pt}
\end{center}

La prossima tappa è Goa Lawah, il tempio che avevamo superato all'andata perché ancora chiuso. Durante il tragitto siamo superati da un macchinone nero con vetri oscurati che viaggia velocemente noncurante degli altri e suona il clacson all'impazzata. Noto che i numeri sulla targa sono tutti nove: DK9999. Penso si tratti di qualche persona di potere che ha acquistato la targa, poiché per alcune religioni di matrice indiana il numero nove è sacro. Mi avevano raccontato che in India certi personaggi pagavano l’equivalente di migliaia di euro per avere un numero di cellulare con più numeri nove possibile.

Arrivati a Goa Lawah, prima di recarci al tempio siamo attratti da una cerimonia indù che si svolge di fronte al mare. Un centinaio di persone — uomini, donne e bambini, tutti vestiti in abiti tradizionali bianchi — portano ceste piene di offerte verso il mare per benedirle e porgere doni. Tra queste offerte vi sono anche animali vivi, tra cui un'anatra che, una volta offerta, viene presa al volo da un ragazzo indonesiano sulla spiaggia con gli amici, che se la porta via come fosse un trofeo. Siamo quasi più attratti dalla cerimonia che dal tempio e il numero di fotografie che scattiamo lo dimostra.

Passiamo poi nell'area limitrofa al tempio, schivando le venditrici ambulanti che qui sono particolarmente assillanti: ci rincorrono e ci toccano, cercando di metterci al collo collane di bambù vendute come "gift" e provando a venderci il Sarong dicendo che è obbligatorio per entrare. Le loro maniere ci fanno allontanare quasi di corsa, senza nemmeno degnarle di uno sguardo.

Prima di entrare facciamo tappa in un negozietto locale — la solita baracca improvvisata che vende prodotti vari — dove compriamo chips e caffè. Le chips non sono di banana come avevamo supposto, ma di patata e pure piccanti, mentre il caffè è servito in un bicchiere di vetro già zuccherato. Siamo molto dubbiosi sull'igiene del bicchiere in generale, ma speriamo che la temperatura lavica del caffè abbia ucciso i batteri di quella brodaglia che, a dispetto del contesto, ha un sapore ottimo: ricorda molto il caffè assaggiato nei villaggi Sasak a Lombok.

All'ingresso del tempio una guida locale ci propone il tour guidato, ma dal momento che il tempio è piccolo e abbiamo un programma da rispettare preferiamo girare da soli, lasciandolo molto deluso — anche se si rifà presto con un'altra coppia di turisti. Alla fine veniamo comunque braccati dalle venditrici che riescono a vestirci coi Sarong e a regalarci delle collanine. Entriamo.

\begin{center}
    % Decorative line
    \noindent\rule{0.6\textwidth}{0.4pt}
\end{center}

\textit{Pura Goa Lawah, in balinese "Tempio della grotta dei pipistrelli", è un tempio indù situato a Klungkung, fondato nell'XI secolo da Mpu Kuturan, uno dei primi sacerdoti che introdussero l'induismo a Bali. È spesso incluso tra i Sad Kahyangan Jagad, i "sei santuari del mondo", i sei luoghi di culto più sacri di Bali.}

\begin{center}
    % Decorative line
    \noindent\rule{0.6\textwidth}{0.4pt}
\end{center}

Si raggiunge quasi immediatamente la grotta dei pipistrelli, che ci lascia a bocca aperta. Elisa, che aveva letto di questa grotta, continua a ripetere con voce atonale: "Non avevo capito. Non avevo capito. Non avevo capito." Il soffitto, che osserviamo da debita distanza, è completamente ricoperto da pipistrelli neri. Saranno letteralmente migliaia. Tanti sono fermi appesi a testa in giù e altri svolazzano avanti e indietro. I loro versi si sentono fin da fuori, insieme all'intenso odore del guano. Abbiamo trovato la caverna di Batman!

All'interno della grotta vi sono aree di preghiera con donne indonesiane che puliscono continuamente e altre che portano doni e offerte. Improvvisamente il gruppo che prima era davanti al mare si riversa all'interno del tempio per una cerimonia: intuiamo si tratti di un rito funebre — tutti vestiti di bianco — e che dopo aver sparso le ceneri del defunto in mare abbiano portato un simbolo del proprio caro alla grotta per pregare affinché possa reincarnarsi bene. Almeno questa è l'idea che ci siamo fatti.

Collezioniamo foto e video finché possiamo, poi i fedeli si moltiplicano e siamo costretti a farci da parte. Osserviamo la cerimonia per un po' mentre Eli scatta soprattutto primi piani, poi proseguiamo. Piccolo ma ricco di fascino, l'aver assistito alla cerimonia gli ha fatto guadagnare parecchi punti. Di questo tempio avevamo letto: "si può potenzialmente anche rimanere delusi, ma bisogna saperne captare l'energia mistica."

Usciti dal tempio restituiamo i Sarong e ovviamente le collanine non si rivelano un regalo ma un amo per farci comprare souvenir: restituiamo tutto e quando non compriamo nulla ci guardano addirittura male. E sticazzi, aggiungerei, ma non saprei come dirlo in balinese.

Prima di raggiungere la macchina ci fermiamo in un ampio spiazzo dove pranziamo divorando i panini di Eli, che sono incredibilmente buoni nonostante gli ingredienti di dubbia origine. Abbiamo anche una bottiglia di Coca-Cola che è rimasta tutto il tempo nello zaino e che Elisa apre con la sua proverbiale delicatezza, facendone schizzare più di metà sulle gambe, che resteranno appiccicose per il resto della giornata.

Arrivati alla macchina troviamo il nostro autista che ci aspetta e mi decido finalmente a chiedergli come si chiami: è Ketut! Provo a schivare la figuraccia dicendo che mi ero solo dimenticato il nome, ma serve a poco — e d'altronde temo che ci siano abituati.

\begin{center}
    % Decorative line
    \noindent\rule{0.6\textwidth}{0.4pt}
\end{center}

L'itinerario riprende e, complici la stanchezza accumulata, la pancia piena e Elisa che inizia a leggermi il significato dei simboli sacri osservati nei templi, sono io stavolta a crollare addormentato quasi subito, per risvegliarmi direttamente all'arrivo al nuovo tempio. Ricordo solo che aveva cominciato a spiegarmi il significato della svastica nella cultura indù, quindi è giusto darle un po' di spazio almeno qui.

\begin{center}
    % Decorative line
    \noindent\rule{0.6\textwidth}{0.4pt}
\end{center}

\textit{La svastica rivolta verso destra simboleggia l'evoluzione del primo principio della creazione, l'evoluzione dell'innocenza e della divinità. Quella rivolta a sinistra è simbolo di distruzione — ma della distruzione degli ostacoli alla via dell'evoluzione divina. Le due direzioni sono complementari e la svastica, esprimendole entrambe, è il simbolo dell'"autosviluppo della santità che distrugge tutti gli ostacoli nel processo di esecuzione del piano divino", nonché il "simbolo della natura divina dell'uomo".}

\begin{center}
    % Decorative line
    \noindent\rule{0.6\textwidth}{0.4pt}
\end{center}

Credo che questo non lasci dubbi sul perché mi sia addormentato.

Anche al tempio di Gunung Kawi Tampaksiring ci prestano dei Sarong insieme ai biglietti d'ingresso. Iniziamo a scendere per una stradina che offre una meravigliosa vista delle risaie e passa tra le botteghe dei venditori, meno pressanti dei precedenti ma più furbi. In tanti ci propongono oggetti come Sarong o souvenir di legno intagliato e quando ce li mettono sotto il naso esclamano "un euro" o "un dollar." I turisti più distratti ci cascano e dopo essersi fermati a vedere e provare i prodotti si sentono rivelare che quello che avrebbero detto è "in dollar" — ossia che è permesso pagare in dollari.

\begin{center}
    % Decorative line
    \noindent\rule{0.6\textwidth}{0.4pt}
\end{center}

\textit{Il tempio Gunung Kawi, popolarmente noto come la Valle dei Re Balinesi, è un tempio e complesso funerario dell'XI secolo a Tampaksiring, che si estende su entrambi i lati del fiume Pakerisan. È composto da dieci candi (santuari) scavati nella roccia, scolpiti in nicchie alte circa sette metri della parete scoscesa della scogliera. Si pensa che questi monumenti funebri siano dedicati al re Anak Wungsu della dinastia Udayana e alle sue regine. Sul santuario settentrionale è leggibile un'iscrizione: "Haji Lumahing Jalu", che significa "il re costruì un tempio qui".}

\begin{center}
    % Decorative line
    \noindent\rule{0.6\textwidth}{0.4pt}
\end{center}

Ci sono pochi turisti e possiamo esplorarlo abbastanza in autonomia: scattiamo foto su un bellissimo ponte di pietra in mezzo alla folta vegetazione, visitiamo il tempio di preghiera interno, contempliamo gli enormi santuari nella roccia e osserviamo i locali mentre intagliano a mano noci di cocco per creare decori ornamentali. Il tempio si trova all'interno di una foresta ricca di risaie, sul fondo di un avvallamento, e tutto rende ancora più suggestiva la visita.

Ciò che più ci lascia senza parole — o con parole che è meglio non riportare qui — è un gruppo di ragazze russe che si fotografano in posa davanti al tempio e, dopo una prima sessione di foto, si cambiano d'abito e ricominciano a scattare. Instagram fa male.

La visita è abbastanza veloce ma uscire non è altrettanto facile: dobbiamo risalire tutti i gradini con un'umidità folle. Per nostra fortuna il sole è temporaneamente coperto e questo facilita la salita. A proposito dell'umidità, Elisa ha sviluppato una vera ossessione per le muffe sulle pareti di pietra: ormai le vede dappertutto e le commenta con ardore.

\begin{center}
    % Decorative line
    \noindent\rule{0.6\textwidth}{0.4pt}
\end{center}

All'uscita saliamo in auto in direzione dell'ultimo tempio: Tirta Empul.

\begin{center}
    % Decorative line
    \noindent\rule{0.6\textwidth}{0.4pt}
\end{center}

\textit{Il tempio di Tirta Empul, situato vicino a Tampaksiring, è un tempio indù balinese d'acqua fondato nel 962 d.C. durante la dinastia Warmadewa. Il nome deriva dalla sorgente sotterranea "Tirta Empul", considerata sacra dagli indù balinesi e origine del fiume Pakerisan. Il complesso è noto per la sua struttura per il bagno, dove i fedeli praticano il rituale di purificazione Melukat utilizzando l'acqua sacra della sorgente, chiamata amritha. Su una collina sovrastante, è stata costruita una villa per la visita del presidente Sukarno nel 1954, oggi usata come casa di riposo per ospiti importanti.}

\begin{center}
    % Decorative line
    \noindent\rule{0.6\textwidth}{0.4pt}
\end{center}

Dal numero di macchine nel parcheggio è facile intuire che siamo giunti in una meta molto turistica. Questo tempio è famoso per la possibilità di partecipare al rito di purificazione all'interno della piscina. Oltre alla visita tradizionale è possibile acquistare un pacchetto che comprende offerta agli dei e bagno nella piscina sotto i getti d'acqua provenienti dalle bocche delle statue. Tantissimi occidentali vestono Sarong verdi sintetici e cinture rosse, accompagnati da guide che forniscono indicazioni in maniera superficiale e distratta. All'interno della piscina questi turisti — i gruppi di più giovani in particolare — non fanno altro che ridere e urlare, raramente zittiti in quello che dovrebbe essere un luogo sacro.

In mezzo a questa caciara osserviamo anche chi si è recato al tempio con vero intento religioso, soprattutto indiani — dell'India stavolta — che indossano abiti da cerimonia personali e vivono i momenti di preghiera con reale devozione. In questo momento mi immagino sulle sponde della piscina impugnando un lungo bastone, lo Spiezzaschièn, a elargire bastonate qua e là a chi urla e fa casino, mentre intono a gran voce la parabola di Tatanaele e Isappo e ogni tanto, tra una bastonata e l'altra, sputo sulla testa di un turista francese per mio personale giubilo. Capisco che questa trovata turistica serva a finanziare la cura del tempio, ma sarebbe stato meglio impedire ai non credenti di praticare il rito nelle piscine. Così sembra il parco acquatico di Asti Lido.

Iniziamo a essere stanchi della giornata nonostante i ripetuti sonnellini in auto. Oltre alle piscine il tempio ospita gazebo adibiti alle offerte e alle preghiere — che ormai conosciamo a memoria — così il nostro tour si avvia presto alla fine.

\begin{center}
    % Decorative line
    \noindent\rule{0.6\textwidth}{0.4pt}
\end{center}

Torniamo alla macchina e ripartiamo verso l'ultima meta: le famose risaie di Tegallalang. Anche questa, essendo la più vicina a Ubud, è molto turistica e l'autista ci lascia in una zona di scarico passeggeri per parcheggiare più lontano. L'area di accesso sembra un centro commerciale, con bar, ristoranti e un ampio bancone dove si può scegliere il pacchetto preferito. Oltre all'accesso standard è possibile fotografarsi in vestiti lunghi colorati sulle altalene legate tra le palme che dondolano sul vuoto, fare la zipline sopra le risaie, pedalare su una bicicletta su un filo di metallo teso ad altezza considerevole e chi più ne ha più ne metta. Di nuovo, sembra che Instagram abbia colonizzato i tour.

Acquistiamo un biglietto standard ed entriamo.

\begin{center}
    % Decorative line
    \noindent\rule{0.6\textwidth}{0.4pt}
\end{center}

\textit{Le risaie di Tegallalang, oggi Patrimonio dell'UNESCO, sono famose per l'antico sistema di irrigazione balinese chiamato Subak, tramandato nel corso dei secoli. Le risaie sono costruite su ripidi pendii e offrono un panorama di verde intenso e paesaggi spettacolari.}

\begin{center}
    % Decorative line
    \noindent\rule{0.6\textwidth}{0.4pt}
\end{center}

Qui ci facciamo scattare qualche foto da un simpatico addetto all'interno di un enorme cuore realizzato con piante e fiori intrecciati, che ci regala anche un album fotografico personalizzato. Decidiamo di non attraversare tutte le risaie a piedi perché siamo stanchi morti. In queste vacanze abbiamo visto risaie in lungo e in largo, conosciuto da Oshi la storia della coltivazione del riso e visto come viene conservato negli antichi villaggi Sasak: adesso vogliamo solo goderci la vista panoramica di quelle che sono sicuramente le risaie più suggestive del viaggio.

Guadagniamo così un po' di tempo che usiamo per una merenda: involtini vietnamiti per me e un piatto di wrap indiani per Elisa, che tenta in tutti i modi di non averli piccanti ma anche questa volta non riesce a essere soddisfatta. L'attesa è un po' lunga, ma d'altronde qui preparano tutto da zero nonostante la folla.

Terminato lo spuntino usciamo, una navetta ci accompagna alla macchina e con Ketut ripartiamo verso casa. Parlando del più e del meno Ketut si offre di accompagnarci a Kuta due giorni dopo, per la tappa finale del viaggio. Lo prendiamo seriamente in considerazione: per trasporti lunghi tra città difficilmente avremmo potuto usare le app di car sharing.

\begin{center}
    % Decorative line
    \noindent\rule{0.6\textwidth}{0.4pt}
\end{center}

Giunto a Ubud il traffico è come al solito congestionato e chiediamo a Ketut di lasciarci in prossimità dell'albergo per percorrere l'ultimo tratto a piedi più velocemente. Il tour in totale è durato dodici ore e siamo stanchissimi, ma lo spuntino alle risaie non ci ha saziati. Dopo una veloce doccia siamo nuovamente in giro alla ricerca di un posto per cenare.

Dopo la sconfitta della sera precedente lascio carta bianca a Elisa e così attraversiamo mezza Ubud per ritrovarci a cenare nel ristorante di un albergo: una bellissima location in mezzo al verde, seduti attorno a un tavolino su un patio di bambù, davanti a un menù che abbiamo sfogliato con cura prima di sederci ma che alla fine "non mi piace niente." La dinamica ancora ora mi è poco chiara.

Io provo per la prima volta degli ottimi Chicken Satay — spiedini di pollo con salsa di arachidi e salsa di soia agrodolce — mentre Eli divora un Nasi Goreng con i suoi ormai immancabili involtini primavera. Per dolce, una crepe verde al cocco e banana. Quasi tutti i miei pasti sono accompagnati dalla Bintang, mentre Elisa prende la versione al limone, la Bintang Radler, che assomiglia tanto a una panaché.

Terminata la cena, foto di rito alla location e ci avviamo verso il Gayatri. Una volta in camera ordiniamo la colazione per l'indomani e intanto Elisa mi legge dal telefono un po' della storia dei templi visitati oggi. Crollo dopo neanche un minuto e per la prima sera non riesco a scrivere nemmeno una parola di diario. So già che dovrò recuperare tutto durante l'interminabile viaggio di ritorno.

\pagestyle{empty} % disable numbers

%coppia
\splitImageLeft{images/day13/PXL_20240821_001813002.jpg}
\splitImageRight{images/day13/PXL_20240821_001813002.jpg}
%

%coppia
\portraitFullscreen{images/day13/PXL_20240821_003521548.PORTRAIT.ORIGINAL.jpg}
\twoImagesLargeMargins{images/day13/PXL_20240821_003222037.MP.jpg}
                      {images/day13/PXL_20240821_004613667.jpg}
%

%coppia
\landscapeThinMargin{images/day13/PXL_20240821_005621826.jpg}
\twoImagesLargeMargins{images/day13/PXL_20240821_010356405.MP.jpg}
                      {images/day13/PXL_20240821_005328791.MP.jpg}
%

%coppia
\splitImageLeft{images/day13/PXL_20240821_010916028.MP.jpg}
\splitImageRight{images/day13/PXL_20240821_010916028.MP.jpg}
%

%coppia
\landscapeThinMargin{images/day13/PXL_20240821_005915545.jpg}
\twoImagesLargeMargins{images/day13/PXL_20240821_010055057.MP.jpg}
                      {images/day13/PXL_20240821_010040656.jpg}
%

%coppia
\splitImageLeft{images/day13/PXL_20240821_011017456.jpg}
\splitImageRight{images/day13/PXL_20240821_011017456.jpg}

\twoImagesLargeMargins{images/day13/PXL_20240821_011901810.jpg}
                      {images/day13/PXL_20240821_011937910.MP.jpg}
\portraitFullscreen{images/day13/PXL_20240821_012052725.jpg}
%

%coppia
\twoImagesLargeMargins{images/day13/PXL_20240821_011413373.jpg}
                      {images/day13/PXL_20240821_011426681.MP.jpg}
\twoImagesLargeMargins{images/day13/PXL_20240821_025603218.jpg}
                      {images/day13/PXL_20240821_025607256.jpg}
%

%coppia
\splitImageLeft{images/day13/PXL_20240821_025623570.jpg}
\splitImageRight{images/day13/PXL_20240821_025623570.jpg}
%

%coppia
\splitImageLeft{images/day13/PXL_20240821_025507333.MP.jpg}
\splitImageRight{images/day13/PXL_20240821_025507333.MP.jpg}
%

%coppia
\landscapeThinMargin{images/day13/PXL_20240821_032006558.jpg}
\twoImagesLargeMargins{images/day13/PXL_20240821_032144993.jpg}
                      {images/day13/PXL_20240821_032501810.jpg}
%

%coppia
\landscapeThinMarginCentered{images/day13/PXL_20240821_033139413.MP.jpg}
\landscapeThinMargin{images/day13/PXL_20240821_033256541.jpg}
%

%coppia
\twoImagesLargeMargins{images/day13/PXL_20240821_033245775.jpg}
                      {images/day13/PXL_20240821_032929504.jpg}
\landscapeThinMargin{images/day13/PXL_20240821_033256541.jpg}
%

%coppia
\landscapeThinMargin{images/day13/PXL_20240821_034055227.jpg}
\twoImagesLargeMargins{images/day13/PXL_20240821_034228822.MP.jpg}
                      {images/day13/PXL_20240821_034554519.MP.jpg}
%

%coppia
\landscapeThinMargin{images/day13/PXL_20240821_052428498.jpg}
\twoImagesLargeMargins{images/day13/PXL_20240821_052642580.jpg}
                      {images/day13/PXL_20240821_052539189.jpg}
%

%coppia
\twoImagesLargeMargins{images/day13/PXL_20240821_055058735.jpg}
                      {images/day13/PXL_20240821_054639237.jpg}
\portraitFullscreen{images/day13/PXL_20240821_053410031.jpg}
%

%coppia
\threeImagesThinMargins{images/day13/PXL_20240821_055327142.MP.jpg}
                       {images/day13/PXL_20240821_060006157.jpg}
                       {images/day13/PXL_20240821_060326584.jpg}
\landscapeThinMarginCentered{images/day13/PXL_20240821_061512226.jpg}
%

%coppia
\splitImageLeft{images/day13/PXL_20240821_061424087.jpg}
\splitImageRight{images/day13/PXL_20240821_061424087.jpg}
%


%coppia
\twoImagesLargeMargins{images/day13/PXL_20240821_063847442.MP.jpg}
                      {images/day13/PXL_20240821_064014192.PORTRAIT.jpg}
\twoImagesLargeMargins{images/day13/PXL_20240821_064422695.jpg}
                      {images/day13/PXL_20240821_064538256.MP.jpg}
%

%coppia
\splitImageLeft{images/day13/PXL_20240821_063855238.jpg}
\splitImageRight{images/day13/PXL_20240821_063855238.jpg}
%

%coppia
\landscapeThinMargin{images/day13/PXL_20240821_074415757.MP.jpg}
\landscapeThinMargin{images/day13/PXL_20240821_070402689.jpg}
%

%coppia
\portraitFullscreen{images/day13/PXL_20240821_075936482.PORTRAIT.jpg}
\threeImagesThinMargins{images/day13/PXL_20240821_080559627.MP.jpg}
                       {images/day13/PXL_20240821_080903116.jpg}
                       {images/day13/PXL_20240821_084155197.MP.jpg}
%

%coppia
\landscapeThinMargin{images/day13/PXL_20240821_130857922.MP.jpg}
\landscapeThinMargin{images/day13/PXL_20240821_134939081.NIGHT.jpg}
\landscapeThinMargin{images/day13/PXL_20240821_135832211.NIGHT.jpg}
\landscapeThinMargin{images/day13/PXL_20240821_135908402.NIGHT.jpg}
%

\landscapeLargeMarginCentered{images/day13/PXL_20240821_130841328.jpg}

\pagestyle{fancy} % enable numbers